\chapter*{Preface}

Toss a fair coin a million times. The law of large numbers tells us that the proportion of heads will probably be close to one half. The central limit theorem describes the remaining error: its standard deviation is $1/2000$, and its distribution is approximately Gaussian. But how accurate is that approximation? What happens when the variance is infinite? How unlikely is it to see sixty percent heads? And what can we say about an entire random path, rather than its endpoint?

These questions lead into the first part of this book. We begin with zero--one laws and random series, then prove the strong law and Kolmogorov's three-series theorem. Donsker's invariance principle shows that a centered walk with independent identically distributed increments of finite positive variance, suitably rescaled, converges in distribution to Brownian motion. The law of the iterated logarithm describes its almost-sure fluctuations on a different scale. Finally, the Lindeberg--Feller theorem and Berry--Esseen inequality tell us when Gaussian approximation is valid and how large its error can be.

The second part asks what replaces the Gaussian law when its hypotheses fail, and how to measure probability more precisely when they hold. Infinitely divisible distributions arise as limits of sums of infinitesimal independent arrays; stable laws describe limits of suitably normalized sums of identically distributed variables. Stein's method gives quantitative normal and Poisson approximations. Cram\'er's and Sanov's theorems describe rare events on an exponential scale, while concentration inequalities bound the fluctuations of functions of many independent inputs. The examples range from the longest increasing subsequence of a permutation to the preservation of distances under random projection.

In the third part, independence is no longer the organizing principle. Martingales describe information revealed over time. Exchangeability leads to de Finetti's theorem, and infinite products of probability measures lead to Kakutani's dichotomy. Ergodic theory replaces independence by invariance: Birkhoff's theorem explains continued-fraction digit frequencies, and Furstenberg's argument proves the existence of three-term arithmetic progressions in sets of positive density. Mod-$\varphi$ convergence and the Poisson--Dirichlet law then bring probability into contact with arithmetic and combinatorics.

The last two chapters return to concrete objects. How many transpositions, or how many riffle shuffles, make a large deck nearly random? Both models exhibit cutoff, a sharp transition on their respective mixing-time scales. A final collection of problems brings together the tools of the volume, including permutations, random series, random polynomials, Minkowski's question mark and the orbit of a Cantor point.

The aim is not merely to recognize the names of the theorems, but to understand their assumptions, proofs and limits. Results used from earlier volumes are identified where they enter. Some deeper theorems are quoted rather than proved, including the infinitesimal-array converse and nonsymmetric stable-law classification in Chapter~5, and the Selberg--Delange formula in Chapter~11; the last is proved in Volume~15. Other quoted results are marked in the text. Each chapter includes graded exercises or problems, with hints and worked solutions. The figures are intended to make mathematical mechanisms visible, not to substitute a numerical experiment for a proof.

This volume assumes the probability of Volume~4, the measure and integration of Volume~7 and the Hilbert spaces of Volume~8. A reader may begin after Volume~8, leaving the representation-theoretic portions of Chapter~13 until after Volume~10. Arithmetic references, including Erd\H{o}s--Kac in Chapter~11, are identified in context.

The governing question of \emph{Order and Chance} is how regularity emerges from apparent randomness. Here the answer takes many forms: a limiting distribution, a concentration inequality, an invariant measure, a recurrence theorem or a mixing-time estimate. The next volume studies the geometry of the continuous random paths that have appeared here.
